\documentclass[11pt]{article}
\usepackage[margin=1in]{geometry}
\usepackage{booktabs}
\usepackage{amsmath}
\usepackage{listings}
\usepackage{xcolor}
\usepackage{float}
\usepackage{hyperref}

\lstset{
  basicstyle=\ttfamily\small,
  keywordstyle=\color{blue},
  commentstyle=\color{gray},
  breaklines=true,
  frame=single,
  language=Python,
}

\title{Triton Tensor Product Kernel Optimization:\\
  Benchmarking \boldmath$C = A_r(A_s(A_t B))$ on NVIDIA H100}
\author{Dylan Jude}
\date{March 2026}

\begin{document}
\maketitle

\section{Introduction}

This report documents the development and optimization of a Triton
implementation of the 3D tensor product
\[
  C = A_r \bigl(A_s \bigl(A_t \, B\bigr)\bigr),
\]
used in spectral element methods. The operator matrices
$A_r, A_s, A_t \in \mathbb{R}^{M \times K}$ are applied successively to
a batch of $N$ input tensors $B \in \mathbb{R}^{N \times K \times K \times K}$
to produce $C \in \mathbb{R}^{N \times M \times M \times M}$.

The computation is \emph{sum-factorized} into three stages:
\begin{align}
  T_1[n,i,j,m] &= \sum_{k} A_t[m,k] \; B[n,i,j,k],    \label{eq:s1} \\
  T_2[n,i,b,c] &= \sum_{j} A_s[b,j] \; T_1[n,i,j,c],   \label{eq:s2} \\
  C[n,a,b,c]   &= \sum_{i} A_r[a,i] \; T_2[n,i,b,c].   \label{eq:s3}
\end{align}
Sum-factorization reduces the per-element FLOP count from $O(K^3)$ for a
dense matrix--vector product to $O(3K)$, at the cost of intermediate
storage for $T_1$ and $T_2$.

All benchmarks use $M=8$, $K=4$, $N=65{,}536$ in double precision
(float64) on an NVIDIA H100 80GB HBM3 (SM~9.0).

\subsection{FLOP counting}

Each stage performs a dot product of length $K$ for every output element:
\begin{itemize}
  \item Stage~1: $K^2 M$ outputs $\times$ $2K$ FLOPs $= 2K^3 M$
  \item Stage~2: $K M^2$ outputs $\times$ $2K$ FLOPs $= 2K^2 M^2$
  \item Stage~3: $M^3$ outputs $\times$ $2K$ FLOPs $= 2K M^3$
\end{itemize}
Total per slice: $2K(K^2 M + K M^2 + M^3)$.  For the full batch:
\[
  \text{FLOP} = N \cdot 2K \cdot (K^2 M + K M^2 + M^3).
\]
With $M=8$, $K=4$, $N=65{,}536$: $\text{FLOP} = 469{,}762{,}048 \approx 0.47\,\text{GFLOP}$.

\textbf{Note:} The original C++ code contained a bug in the FLOP formula,
using $(2N-1)(K^2 M + K M^2 + M^3)$ instead of the correct expression above.
This was fixed as part of this work.

\section{Existing implementations}

The project already contained C++/OCCA GPU kernels and a cuBLAS GEMM
baseline.  Benchmark results on H100:

\begin{table}[H]
\centering
\begin{tabular}{lrr}
\toprule
\textbf{Method} & \textbf{Time (s)} & \textbf{TFLOPS} \\
\midrule
cuBLAS GEMM (dense, no sum-factorization) & $1.65 \times 10^{-4}$ & 26.1 \\
OCCA fused kernel (sum-factorized, shared mem) & $1.12 \times 10^{-4}$ & 4.18 \\
C++ OpenMP (CPU) & $6.70 \times 10^{-3}$ & --- \\
\bottomrule
\end{tabular}
\caption{Pre-existing implementations.  cuBLAS reports higher TFLOPS because
it performs $\sim\!9\times$ more arithmetic (dense GEMM of $M^3 \times K^3$);
the OCCA kernel is faster in wall-clock time.}
\label{tab:existing}
\end{table}

\section{Triton kernel development}

\subsection{Variant 1: Three separate stage kernels}

The initial Triton implementation uses three separate kernel launches,
one per contraction stage~\eqref{eq:s1}--\eqref{eq:s3}.  Each kernel
launches $N$ programs (one per batch element), with each program
vectorizing over the output elements of that stage using
\texttt{tl.arange}.  The contraction index is iterated via
\texttt{tl.static\_range} for compile-time unrolling.

Intermediates $T_1$ and $T_2$ are allocated as global-memory tensors
and passed between kernel launches.

\begin{table}[H]
\centering
\begin{tabular}{lrrr}
\toprule
\textbf{Stage} & \textbf{Output elements/slice} & \textbf{BLOCK size} & \textbf{Contraction iters} \\
\midrule
Stage 1 & $K^2 M = 128$  & 128 & $K = 4$ \\
Stage 2 & $K M^2 = 256$  & 256 & $K = 4$ \\
Stage 3 & $M^3 = 512$    & 512 & $K = 4$ \\
\bottomrule
\end{tabular}
\end{table}

\textbf{Memory traffic per slice} (read + write):
$K^3 + 2 K^2 M + 2 K M^2 + M^3 = 1344$ doubles $= 10{,}752$ bytes.
\textbf{Arithmetic intensity}: $7{,}168 / 10{,}752 = 0.67$ FLOP/byte --- deeply
bandwidth-bound on H100 (crossover $\approx 10$ FLOP/byte).

\begin{center}
\fbox{%
\begin{minipage}{0.7\textwidth}
\centering
\textbf{Result:} $\mathbf{2.55 \times 10^{-4}}$~\textbf{s} \quad (1.84 TFLOPS)
\end{minipage}}
\end{center}

\subsection{Variant 2: Fused kernel with M\texorpdfstring{$^2$}{2} vectorization}

The first attempt at fusion vectorized over the $(b,c)$ output plane
($M^2 = 64$ lanes) and looped over the output index $a$ and all
contraction indices $(i,j,k)$ via nested \texttt{tl.static\_range}.
Stage~1 was recomputed $M$ times (once per~$a$), trading extra FLOPs
for eliminated intermediate memory traffic.

This performed very poorly: \textbf{1.31 ms} (0.36 TFLOPS).  The reasons:
\begin{itemize}
  \item Only 64 lanes $=$ 2 warps per program --- poor GPU occupancy.
  \item $M \times K \times K \times K = 672$ unrolled loop iterations
    generated an enormous instruction footprint, likely thrashing the
    instruction cache.
\end{itemize}

\subsection{Variant 3: Fused kernel with M\texorpdfstring{$^3$}{3} vectorization}

To address the occupancy problem, the fused kernel was rewritten to
vectorize over all $M^3 = 512$ output elements $(a,b,c)$, giving 16
warps per program.  The innermost $k$ contraction was partially
sum-factorized:

\begin{lstlisting}
for i in tl.static_range(K):
    ar_val = tl.load(Ar_ptr + a * K + i, mask=mask)
    for j in tl.static_range(K):
        as_val = tl.load(As_ptr + b * K + j, mask=mask)
        ar_as = ar_val * as_val
        tmp = tl.zeros([BLOCK], dtype=tl.float64)
        for k in tl.static_range(K):
            at_val = tl.load(At_ptr + c * K + k, mask=mask)
            b_val = tl.load(B_ptr + n*K*K*K + i*K*K + j*K + k)
            tmp += at_val * b_val
        acc += ar_as * tmp
\end{lstlisting}

Loop iterations reduced from 672 to $K^2 \times K = 64$.
Operator matrices ($3 \times 32$ doubles) stay in L1 cache; $B$ values
are scalar broadcasts.

\begin{center}
\fbox{%
\begin{minipage}{0.7\textwidth}
\centering
\textbf{Result:} $\mathbf{2.89 \times 10^{-4}}$~\textbf{s} \quad (1.63 TFLOPS)
\end{minipage}}
\end{center}

A $4.5\times$ improvement over variant~2, but still slightly slower than the
3-stage separate kernels.  The extra FLOPs from recomputing stage~1
$M=8$ times are not offset by the eliminated memory traffic, because
H100's 3.35~TB/s bandwidth makes the intermediates relatively cheap.

\subsection{Variant 4: Fused 3-stage kernel (L1-cached intermediates)}

The final attempt ran all three contraction stages sequentially within
a single Triton program, writing intermediates to global-memory scratch
buffers.  The hypothesis was that since the same program writes then
immediately reads back, data would stay in L1 cache, giving: single
kernel launch, true sum-factorization, and fast intermediate access.

All stages shared the same \texttt{BLOCK}~$= 512$ (the maximum),
with masking for stages 1 ($n_\text{out} = 128$) and 2
($n_\text{out} = 256$).

\begin{center}
\fbox{%
\begin{minipage}{0.7\textwidth}
\centering
\textbf{Result:} $\mathbf{4.34 \times 10^{-4}}$~\textbf{s} \quad (1.08 TFLOPS)
\end{minipage}}
\end{center}

This was the slowest Triton variant.  The sequential execution of three
stages within one program means:
\begin{itemize}
  \item Stages 1 and 2 waste 75\% and 50\% of lanes respectively.
  \item Programs take $\sim\!3\times$ longer, reducing concurrent occupancy.
  \item Cross-lane visibility of intermediates through global memory
    requires implicit barriers that add latency.
\end{itemize}

\section{Summary of results}

\begin{table}[H]
\centering
\begin{tabular}{llrrl}
\toprule
\textbf{Kernel} & \textbf{Framework} & \textbf{Time (ms)} & \textbf{TFLOPS} & \textbf{Bound} \\
\midrule
OCCA fused                      & OCCA/CUDA & 0.112 & 4.18 & Bandwidth \\
Triton 3-stage                  & Triton    & 0.255 & 1.84 & Bandwidth \\
Triton fused ($M^3$ vec.)       & Triton    & 0.289 & 1.63 & Compute \\
Triton fused 3-stage            & Triton    & 0.434 & 1.08 & Occupancy \\
torch.einsum                    & PyTorch   & 0.844 & 0.56 & Overhead \\
Triton fused ($M^2$ vec.)       & Triton    & 1.310 & 0.36 & I-cache \\
\bottomrule
\end{tabular}
\caption{All kernel variants on H100, $M=8$, $K=4$, $N=65{,}536$, float64.
  All TFLOPS values use the same sum-factorized FLOP count for fair comparison.}
\label{tab:summary}
\end{table}

\section{Conclusion}

The best-performing Triton kernel (3-stage, separate launches) achieves
0.255~ms, which is $2.3\times$ slower than the OCCA fused kernel
(0.112~ms).  This gap is \textbf{fundamental to the programming models},
not a tuning issue.

The OCCA kernel succeeds because it directly maps the algorithm's needs:
\begin{enumerate}
  \item \textbf{Shared memory for intermediates.}  $T_1$ and $T_2$ are
    small per-element arrays ($128$--$256$ doubles) that fit in shared
    memory.  By keeping them on-chip, the OCCA kernel eliminates all
    intermediate global memory traffic and achieves true sum-factorization
    with minimal bandwidth cost.
  \item \textbf{Explicit thread synchronization.}  Between stages, OCCA
    emits \texttt{\_\_syncthreads()} barriers so that all threads in a
    block cooperatively produce the intermediate and then consume it.
    This is the key mechanism that Triton lacks.
  \item \textbf{3D thread blocks.}  The OCCA kernel maps the $(i,j,k)$
    spatial indices directly to CUDA thread dimensions with \texttt{@inner}
    loops, giving natural data reuse across the block.
\end{enumerate}

Triton's programming model, by design, abstracts away shared memory and
intra-block synchronization.  This is a strength for many workloads
(matrix multiplication, attention, etc.) but a limitation for
\emph{multi-stage} algorithms where small intermediates must be produced
and consumed cooperatively within a thread block.  Triton programs
vectorize over a 1D \texttt{tl.arange}, which maps poorly to the
changing tensor shapes between stages of sum-factorization.

Despite this limitation, Triton still provides a $3.3\times$ speedup
over \texttt{torch.einsum} and offers a clean, readable Python
implementation that validates correctly against the CPU reference
(max $|\Delta| < 10^{-12}$).  For users without access to OCCA or the
expertise to write CUDA kernels, Triton provides a practical
middle-ground implementation.

\end{document}
